% !TEX root = ../../../main.tex
% 来源：中学数学实验教材·第二册（几何）
% 章节：轨迹与作图 / 轨迹
% 原始文件：5.tex，行 254-277
% 类型：tikz
% ---
\begin{tikzpicture}[scale=1.2]
\tkzDefPoints{0/0/O, 2.5/0/O', 5/0/P}
\draw[dashed] (P)--(-2,0)node[left]{$E'$};
\draw (0,0) circle (2);
\draw[dashed] (O') circle (2.5);
\tkzDefTangent[from with R=P](O, 2)
\tkzGetPoints{A}{B}
\tkzDrawSegments[add=0 and .2](P,A  P,B)
\node at (-.25,-.25){$O$};\node at (.25+2.5,.25){$O'$};
\node at (2.2,0)[above]{$E$};
\tkzDrawPoints(O',O)
\tkzDefShiftPoint[O'](150:2.5){M}
\tkzDefShiftPoint[O'](-142:2.5){N}
\tkzDefShiftPoint[O'](-155:2.5){X}
\tkzAutoLabelPoints[center=O'](M,N,X,A,B,P)
\tkzInterLC(P,M)(O,A) \tkzGetPoints{C'}{C}
\tkzInterLC(P,X)(O,A)\tkzGetPoints{D}{D'}
\tkzInterLC(P,N)(O,A)\tkzGetPoints{F}{F'}
\tkzAutoLabelPoints[center=O](D,F,C,C',D',F')
\draw(C')--(P)--(F');
\draw(D')--(P);
\draw[dashed](M)--(O)--(X);
\draw[dashed](A)--(O);
\end{tikzpicture}
